% The TARA specification, typeset from the Sail model.
%
% `sail --latex --latex-prefix tara` writes sail_latex/commands.tex, one macro per definition of
% the model (\tarafclADDexecute, \taravalstep, ...). This file arranges those macros into a
% reference: the listings are the normative text, the prose around them explains and indexes it.
\documentclass[a4paper,10pt]{article}

\usepackage[T1]{fontenc}
\usepackage{lmodern}
\usepackage{array}
\usepackage{etoolbox}
\usepackage{booktabs}
\usepackage{fancyhdr}
% xcolor, fullpage, listings, hyperref and the Sail listing language, from Sail's own sail.tex.
\input{sail-preamble}
\usepackage[a4paper,left=22mm,right=22mm,top=23mm,bottom=22mm,headheight=14pt,headsep=6mm]{geometry}

% ---------------------------------------------------------------------------------------------
% Page, links and listings
% ---------------------------------------------------------------------------------------------

% The PDF carries no dates, file ID or engine banner, so equal sources give equal bytes.
\pdfinfoomitdate=1
\pdfsuppressptexinfo=-1
\pdftrailerid{}

\hypersetup{
  urlcolor=MediumBlue,
  pdftitle={The TARA Instruction Set Architecture},
  pdfsubject={Reference specification typeset from the Sail model},
  bookmarksdepth=3,
  bookmarksnumbered=true,
  pdfpagemode=UseOutlines,
}
\renewcommand{\sectionautorefname}{\S}
\renewcommand{\subsectionautorefname}{\S}
\renewcommand{\subsubsectionautorefname}{\S}
\renewcommand{\appendixautorefname}{Appendix}

\emergencystretch=2em
\widowpenalty=10000
\clubpenalty=10000
\setlength{\parskip}{0.4\baselineskip plus 0.1\baselineskip}
\setlength{\parindent}{0pt}
\setcounter{tocdepth}{2}
\makeatletter
\def\@listI{\leftmargin\leftmargini \parsep 0pt \topsep 4pt plus 2pt minus 2pt
  \itemsep 2pt plus 1pt minus 1pt}
% Less space around the headings: the instructions are many and close together.
\renewcommand{\section}{\@startsection{section}{1}{\z@}%
  {-3ex \@plus -1ex \@minus -0.2ex}{1.6ex \@plus 0.2ex}{\normalfont\Large\bfseries}}
\renewcommand{\subsection}{\@startsection{subsection}{2}{\z@}%
  {-2.6ex \@plus -0.8ex \@minus -0.2ex}{1ex \@plus 0.2ex}{\normalfont\large\bfseries}}
\renewcommand{\subsubsection}{\@startsection{subsubsection}{3}{\z@}%
  {-2.4ex \@plus -0.6ex \@minus -0.2ex}{0.8ex \@plus 0.2ex}{\normalfont\normalsize\bfseries}}
\makeatother

\pagestyle{fancy}
\fancyhf{}
\fancyhead[L]{\small The TARA Instruction Set Architecture}
\fancyhead[R]{\small\nouppercase{\leftmark}}
\fancyfoot[C]{\small\thepage}
\renewcommand{\headrulewidth}{0.3pt}
\fancypagestyle{plain}{\fancyhf{}\fancyfoot[C]{\small\thepage}%
  \renewcommand{\headrulewidth}{0pt}}

% Listings: a shaded block, flush with the text, its contents indented by \listingmargin. Sail
% lines have at most 100 columns, and the check below makes sure that 100 columns of this font
% fit between the margins.
\definecolor{shade}{gray}{0.95}
\newlength{\listingmargin}
\setlength{\listingmargin}{5pt}
\newcommand{\listingfont}{\fontsize{8.5}{10.2}\ttfamily\selectfont}
\lstset{
  basicstyle=\listingfont,
  showstringspaces=false,
  backgroundcolor=\color{shade},
  frame=single,
  framerule=0pt,
  framesep=2pt,
  framexleftmargin=\listingmargin,
  framexrightmargin=\listingmargin,
  rulecolor=\color{shade},
  xleftmargin=\listingmargin,
  xrightmargin=\listingmargin,
  aboveskip=0pt,
  belowskip=0pt,
}
\newlength{\hundredcolumns}
\settowidth{\hundredcolumns}{\listingfont
  0123456789012345678901234567890123456789012345678901234567890123456789%
  01234567890123456789012345678901}
\ifdim\hundredcolumns>\dimexpr\textwidth-2\listingmargin\relax
  \PackageError{tara}{100 columns of Sail do not fit the page}{Shrink the listing font.}
\fi

% The generated listings link every call to its definition. Definitions that are not typeset
% in this document (Sail's library) have no target: render those calls as plain text.
\let\tarahyperref\hyperref
\renewcommand{\hyperref}[2][]{\ifcsname r@#1\endcsname\tarahyperref[#1]{#2}\else#2\fi}

% Keep a group of listings on one page, and with the paragraph that introduces it.
\newcommand{\keep}[1]{\par\nopagebreak\vspace{3pt}\begin{minipage}{\linewidth}%
  \setlength{\parskip}{0pt}#1\end{minipage}\par\medskip}

% A centred table or diagram, kept with the paragraph that introduces it (a list begins with a
% penalty that allows a break, which \@beginparpenalty removes for this one).
\makeatletter
\newenvironment{display}{%
  \par\nopagebreak\@beginparpenalty=\@M\begin{center}\@beginparpenalty=-51\relax}{\end{center}}
\makeatother

% A link to a label; unlike the generated ones, a typo shows up as an undefined reference.
\newcommand{\xref}[2]{\tarahyperref[#1]{#2}}

% Inline code; a link to a listing (the label is the one the generated macro defines); a
% mnemonic that links to its entry in Section 4.
\newcommand{\code}[1]{\texttt{#1}}
\newcommand{\sym}[2]{\xref{#1}{\code{#2}}}
\newcommand{\mnem}[1]{\xref{taraz#1}{#1}}

\input{sail_latex/commands}

% ---------------------------------------------------------------------------------------------
% The opcode table: one \defop per instruction, in opcode order.
%   \defop{opcode}{hex}{bits}{mnemonic}{format}{syntax}
% The table in Section 3 and the heading of every instruction in Section 4 read this data. The
% three forms of the opcode must agree (\checkop), and tests/test_doc.py checks the mnemonics and
% formats against TARA Studio's opcode table.
% ---------------------------------------------------------------------------------------------

\newcommand{\binval}[5]{\numexpr16*#1+8*#2+4*#3+2*#4+#5\relax}
\newcommand{\checkop}[4]{%
  \ifnum#1="#2 \else\PackageError{tara}{#4: opcode #1 is not hexadecimal #2}{}\fi
  \ifnum#1=\binval#3 \else\PackageError{tara}{#4: opcode #1 is not binary #3}{}\fi}

\def\oplist{}
\newcommand{\defop}[6]{%
  \checkop{#1}{#2}{#3}{#4}%
  \listgadd{\oplist}{#4}%
  \csgdef{op@#4@dec}{#1}\csgdef{op@#4@hex}{#2}\csgdef{op@#4@bin}{#3}%
  \csgdef{op@#4@fmt}{#5}\csgdef{op@#4@syn}{#6}}

\defop{0}{00}{00000}{NOP}{F0}{NOP}
\defop{1}{01}{00001}{HLT}{F0}{HLT}
\defop{2}{02}{00010}{MOV}{F2}{MOV Rd, Rs}
\defop{3}{03}{00011}{LIL}{F3}{LIL Rd, imm}
\defop{4}{04}{00100}{LIH}{F3}{LIH Rd, imm}
\defop{5}{05}{00101}{LDW}{F4}{LDW Rd, off(Rb)}
\defop{6}{06}{00110}{STW}{F4}{STW Rs, off(Rb)}
\defop{7}{07}{00111}{LDB}{F4}{LDB Rd, off(Rb)}
\defop{8}{08}{01000}{STB}{F4}{STB Rs, off(Rb)}
\defop{9}{09}{01001}{ADD}{F1}{ADD Rd, Rs1, Rs2}
\defop{10}{0A}{01010}{SUB}{F1}{SUB Rd, Rs1, Rs2}
\defop{11}{0B}{01011}{ADDI}{F3}{ADDI Rd, imm}
\defop{12}{0C}{01100}{MUL}{F1}{MUL Rd, Rs1, Rs2}
\defop{13}{0D}{01101}{AND}{F1}{AND Rd, Rs1, Rs2}
\defop{14}{0E}{01110}{OR}{F1}{OR Rd, Rs1, Rs2}
\defop{15}{0F}{01111}{XOR}{F1}{XOR Rd, Rs1, Rs2}
\defop{16}{10}{10000}{NOT}{F2}{NOT Rd, Rs}
\defop{17}{11}{10001}{SHL}{F3}{SHL Rd, shamt}
\defop{18}{12}{10010}{SHR}{F3}{SHR Rd, shamt}
\defop{19}{13}{10011}{SLT}{F1}{SLT Rd, Rs1, Rs2}
\defop{20}{14}{10100}{BZ}{F5}{BZ Rs, off}
\defop{21}{15}{10101}{BN}{F5}{BN Rs, off}
\defop{22}{16}{10110}{JMP}{F6}{JMP off}
\defop{23}{17}{10111}{CALL}{F6}{CALL off}
\defop{24}{18}{11000}{RET}{F0}{RET}
\defop{25}{19}{11001}{PUSH}{F7}{PUSH Rs}
\defop{26}{1A}{11010}{POP}{F7}{POP Rd}

% A row of the opcode table.
\newcommand{\oprow}[1]{%
  \csuse{op@#1@dec} & \code{\csuse{op@#1@hex}} & \code{\csuse{op@#1@bin}} &
  \xref{taraz#1}{\code{#1}} & \code{\csuse{op@#1@syn}} &
  \xref{fmt:\csuse{op@#1@fmt}}{\csuse{op@#1@fmt}} & \pageref{taraz#1} \\}

% A group of instructions: a subsection; its introduction stays with the first entry.
\newif\ifgroupstart
\newcommand{\group}[2]{\subsection{#1}\label{#2}\groupstarttrue}

% An instruction of Section 4: its heading, a line with its syntax, opcode and format, the
% prose given as the second argument and the body (the listings), all on one page.
% The label taraz<MNEMONIC> is the target of the links in the generated listings.
\makeatletter
\newenvironment{instruction}[2]{%
  \ifgroupstart\@secpenalty=\@M\fi % no page break between a group's introduction and its entries
  \subsubsection{\texorpdfstring{\code{#1}}{#1}}\label{taraz#1}%
  \@secpenalty=-300\relax\groupstartfalse
  \begin{minipage}[t]{\linewidth}%
  \textbf{Syntax}\enspace\code{\csuse{op@#1@syn}}\hspace{2em}
  \textbf{Opcode}\enspace\csuse{op@#1@dec} (\code{\csuse{op@#1@bin}})\hspace{2em}
  \textbf{Format}\enspace\xref{fmt:\csuse{op@#1@fmt}}{\csuse{op@#1@fmt}}\par\smallskip
  #2\par\smallskip\setlength{\parskip}{0pt}%
}{%
  \end{minipage}\par\medskip}
\makeatother

% A field of a format diagram: n bits wide; the last field of a row closes the box.
\newlength{\bitwidth}
\setlength{\bitwidth}{5mm}
\newcommand{\fld}[2]{\multicolumn{#1}{|c}{\small#2}}
\newcommand{\fldlast}[2]{\multicolumn{#1}{|c|}{\small#2}}
\newcommand{\bitnumber}[1]{{\scriptsize#1}}

\begin{document}
\thispagestyle{plain}

\begin{center}
  {\LARGE\bfseries The TARA Instruction Set Architecture\par}
  \medskip
  {\large Reference specification, typeset from the Sail model\par}
\end{center}

{\hypersetup{linkcolor=black}\setlength{\parskip}{0pt}\tableofcontents}
\clearpage

% =============================================================================================
\section{Introduction}\label{sec:intro}
% =============================================================================================

TARA (\emph{Teaching Architecture for RISC Assembly}) is the 16-bit teaching processor of the
course CS2300 at IIT Madras, developed and maintained by Ayon Chakraborty. A TARA machine has
eight 16-bit general-purpose registers, a unified 2~KB byte-addressed memory, a memory-mapped
input port with five input lines and a memory-mapped $64\times64$ monochrome framebuffer.
Every instruction is one 16-bit word with a 5-bit opcode; 27 of the 32 opcodes are assigned.

\subsection{The model is the specification}\label{sec:model}

The architecture is specified by a model written in Sail, a language for specifying
instruction-set architectures. This document is typeset from that model: each listing is a
definition taken from the Sail source by \code{sail -{}-latex}, so the listings are the normative
text, and the prose around them explains and indexes them. Where the prose and a listing
disagree, the listing is right.

\begin{display}
\begin{tabular}{@{}ll@{}}
\toprule
File & Contents \\
\midrule
\code{model/machine.sail} & State, memory buses, input port, framebuffer (\autoref{sec:state}) \\
\code{model/tara.sail} & The instructions: operands, encoding, decoding, semantics
  (\autoref{sec:formats}, \autoref{sec:isa}) \\
\code{model/step.sail} & The drivers \sym{tarazstep}{step},
  \sym{tarazrunzyinstruction}{run\_instruction}, \sym{tarazreset}{reset}
  (\autoref{sec:execution}) \\
\code{model/syntax.sail} & Assembly syntax (\autoref{sec:syntax}) \\
\bottomrule
\end{tabular}
\end{display}

The model describes the architecture, not the microarchitecture. A call to \sym{tarazstep}{step}
fetches, decodes and executes one instruction; T-states, the control ROM, memory timing and the
status flags Z, N, C and V are not modelled. The flags are ``observational only'' in the RTL: they
drive LEDs and debug outputs, are not architecturally visible and never control a branch.

\subsection{Sources}\label{sec:sources}

The model follows the RTL description published with TARA Studio: the microarchitecture section
of the TARA Studio manual and the TARA hardware manual. Where the RTL description is silent it
follows the ISA reference table of the Studio manual. Quotations in this document are from these
two manuals.
\begin{itemize}\small
\item TARA Studio manual, \emph{Microarchitecture}:
  \url{https://www.cse.iitm.ac.in/~ayon/courses/CS2300/taramanual/tarastudio/index.html#microarch}
\item TARA CPU on the Basys 3, hardware manual:
  \url{https://www.cse.iitm.ac.in/~ayon/courses/CS2300/taramanual/tarahw/index.html}
\end{itemize}
Where neither source says what happens, the model makes an assumption;
\autoref{app:assumptions} lists them. Where the TARA Studio simulator (taracpu 1.2.2)
disagrees with the RTL description, the model follows the RTL description;
\autoref{app:studio} lists the differences.

\subsection{Reading the listings}\label{sec:reading}

A listing shows a definition as written in the Sail source, with Sail's keyword highlighting.
A call of a function, constructor or register access that is defined in this document is a
link to its definition; calls of Sail's library functions are plain text, since the library is
not part of this document. The notation of the model:

\begin{display}
\small
\begin{tabular}{@{}>{\raggedright\arraybackslash}p{0.27\linewidth}%
  p{\dimexpr0.73\linewidth-2\tabcolsep\relax}@{}}
\toprule
Notation & Meaning \\
\midrule
\code{0b1010}, \code{0x2F} & Bit-vector literals: one bit per binary digit, four per
  hexadecimal digit. \code{bitone} is the single bit 1. \\
\code{bits(n)} & A vector of $n$ bits; bit 0 is the least significant. \\
\code{a @ b} & Concatenation; \code{a} supplies the more significant bits. \\
\code{x[i]}, \code{x[h .. l]} & Bit $i$ of \code{x}; bits $h$ down to $l$, inclusive. \\
\code{unsigned(x)}, \code{signed(x)} & The integer value of a bit-vector read as an unsigned
  number, or as a two's-complement number. \\
\code{get\_slice\_int(n, i, 0)} & The low $n$ bits of the integer \code{i}. \\
\code{+}, \code{-}, \code{\&}, \code{|}, \code{\textasciicircum}, \code{not\_vec} & Addition and
  subtraction modulo $2^n$, bitwise and, or, exclusive or and complement, on bit-vectors of
  equal length. Adding an integer to a bit-vector also wraps modulo $2^n$. \\
\code{sail\_shiftleft}, \code{sail\_shiftright} & Logical shifts by an integer count; the
  vacated bits are zero. \\
\code{Some(x)}, \code{None()} & A value that is present or absent (the type \code{option}). \\
\code{X(r)}, \code{X(r) = v} & Read, write register \code{r} (\autoref{sec:registers}). \\
\code{let}, \code{if}, \code{match} & A local binding, a conditional and a pattern match;
  \code{\{ a; b \}} runs \code{a}, then \code{b}. \\
\code{scattered}, \code{function clause} & A function defined by clauses placed next to the
  definitions they belong to: every instruction adds one clause to \code{encode},
  \code{decode} and \code{execute}. \\
\bottomrule
\end{tabular}
\end{display}

% =============================================================================================
\section{Machine state}\label{sec:state}
% =============================================================================================

The architectural state of a TARA machine is held in the Sail registers declared in
\code{machine.sail}: the register file, PC, the memory and the halt latch. Two more registers
are not architectural state proper: \sym{tararegisterzKEYS}{KEYS} latches the input lines and
\sym{tararegisterznextPC}{nextPC} is scratch through which an instruction proposes its
successor.

\subsection{Types}\label{sec:types}

The sizes of the machine are named once.

\keep{\taratypeword\taratypebyte\taratyperegidx\taratypeaddress\taratyperegfile\taratypememory}

A \sym{taratypezword}{word} is the width of a register, of PC and of an instruction; an
\sym{taratypezaddress}{address} is the 11-bit byte address that selects one of the 2048 bytes of
memory; a \sym{taratypezregidx}{regidx} is a register number. The model binds \code{-} to
bit-vector subtraction and \code{\textasciicircum} to bitwise exclusive or; this is notation only,
the operations come from Sail's library.

% The macro names of overloads carry a letter that counts the earlier overloads of the operator
% in Sail's library (B: the second, C: the third); they change when the library does.
\keep{\taraoverloadBzEightoperatorzZerozDzNine\taraoverloadBzEightoperatorzZerozQzNine}

\subsection{Registers}\label{sec:registers}

The register file \sym{tararegisterzGPR}{GPR} holds R0 to R7, each 16 bits wide. R0 is an ordinary
writable register: nothing is hard-wired to zero. Two registers have a role in the instruction
set: R6 is the link register, written by \mnem{CALL} and read by \mnem{RET}, and R7 is the stack
pointer, used by \mnem{PUSH} and \mnem{POP}; the model names them \sym{taraletzlr}{lr} and
\sym{taraletzsp}{sp}. Registers are read and written by number through \sym{tarazX}{X}.

\keep{\tararegisterGPR\taraletlr\taraletsp\taravalrX\tarafnrX\taravalwX\tarafnwX
  \taraoverloadAX\label{tarazX}}

The program counter \sym{tararegisterzPC}{PC} holds the byte address of the next instruction. It
is 16 bits wide, but only 11 of them are kept: the RTL says ``16-bit \textperiodcentered{} +2 per
fetch \textperiodcentered{} masked to 0x7FF''. The model writes PC only through
\sym{tarazpczymask}{pc\_mask} (and sets it to 0 on reset), so PC wraps around at the end of the
2~KB memory.

\keep{\tararegisterPC\taravalpcMask\tarafnpcMask}

The halt latch \sym{tararegisterzHALTED}{HALTED} is set by \mnem{HLT} and cleared by reset
(\autoref{sec:halting}). The scratch register \sym{tararegisterznextPC}{nextPC} holds the
successor of the current instruction while it executes (\autoref{sec:retire}); it is not
architectural.

\keep{\tararegisterHALTED\tararegisternextPC}

\subsection{Memory and buses}\label{sec:memory}

Memory is an array \sym{tararegisterzMEM}{MEM} of 2048 bytes at the addresses \code{0x000} to
\code{0x7FF}. Code, data, stack, the input port and the framebuffer all live in it.

\keep{\tararegisterMEM}

\begin{display}
\begin{tabular}{@{}ll@{}}
\toprule
Addresses & Contents \\
\midrule
\code{0x000}--\code{0x5FE} & Ordinary memory for program, data and stack (1535 bytes) \\
\code{0x5FF} & Input port (\autoref{sec:input}) \\
\code{0x600}--\code{0x7FF} & Framebuffer, 512 bytes (\autoref{sec:framebuffer}) \\
\bottomrule
\end{tabular}
\end{display}

The Studio manual suggests a conventional layout (code from \code{0x000}, data from
\code{0x400}, a stack below \code{0x5F0}); the architecture enforces none of it.

\paragraph{Byte bus.} \sym{tarazreadzybyte}{read\_byte} and
\sym{tarazwritezybyte}{write\_byte} take a 16-bit address and use
only its low 11 bits, so addresses wrap modulo 2048: an effective address computed as
$0-1$ selects byte \code{0x7FF}, and \code{0xF802} selects byte \code{0x002}.

\keep{\taravalreadByte\tarafnreadByte\taravalwriteByte\tarafnwriteByte}

\paragraph{Word bus.} The RTL gives the rule ``word index = addr \guillemotright{} 1
\textperiodcentered{} lane = addr[0]'' and stores words big-endian: the word at address $2k$
has its high byte in byte $2k$ and its low byte in byte $2k+1$. A word access ignores address
bit 0, so an access to the odd address $2k+1$ is an access to the word at $2k$, and the last
word of memory is the one at \code{0x7FE}. A word access is two byte accesses, so a word access
that covers \code{0x5FF} sees the input port (\autoref{sec:input}). Instruction fetch is a word
read (\autoref{sec:step}); \sym{tarazreadzyword}{read\_word} and
\sym{tarazwritezyword}{write\_word} are the word bus.

\keep{\taravalreadWord\tarafnreadWord\taravalwriteWord\tarafnwriteWord}

\subsection{Input port}\label{sec:input}

``A byte read from \code{0x5FF} returns \{3'b000, QUIT, RIGHT, LEFT, DOWN, UP\}.'' The model
keeps the five input lines in \sym{tararegisterzKEYS}{KEYS}, with UP in bit 0 and QUIT in bit 4,
and \code{read\_byte} returns them, zero-extended, whenever the low 11 bits of the address are
\sym{taraletzinputzyport}{input\_port}.

\keep{\taraletinputPort\tararegisterKEYS}

\begin{display}
\begin{tabular}{@{}cccccc@{}}
\toprule
Bits 7--5 & Bit 4 & Bit 3 & Bit 2 & Bit 1 & Bit 0 \\
\midrule
\code{000} & QUIT & RIGHT & LEFT & DOWN & UP \\
\bottomrule
\end{tabular}
\end{display}

The drivers latch the lines once per step (\autoref{sec:execution}), so all reads of
\code{0x5FF} during one instruction, including the fetch of the instruction itself, see the same
value. How the lines are driven is up to the host; TARA Studio uses the arrow keys or W, A, S,
D for the four directions and Q for QUIT. A store to \code{0x5FF} is not special:
\code{write\_byte} writes the underlying memory byte, which no read returns.

\subsection{Framebuffer}\label{sec:framebuffer}

The 64$\times$64 framebuffer occupies \code{0x600} to \code{0x7FF}: 512 bytes, one bit per pixel,
eight pixels per byte with the most significant bit on the left. Row 0 is the bottom of the
display. The RTL gives the mapping: ``Pixel $(x, y)$: byte = \code{0x0600} + $y\times8$ +
$x\div8$, bit = $7 - (x \bmod 8)$''. The model has no separate display state:
\sym{tarazpixel}{pixel} reads the memory and tells whether the pixel $(x, y)$ is lit.

\keep{\taravalpixel\tarafnpixel}

% =============================================================================================
\section{Instruction formats and opcodes}\label{sec:formats}
% =============================================================================================

\subsection{Formats}\label{sec:format-diagram}

An instruction is one 16-bit word, stored big-endian like every word. Bits 15 to 11 hold the
opcode; the other 11 bits are divided into fields by one of eight formats, F0 to F7 (the names
are those of the Studio manual). A field written as zeros is padding: \code{decode} ignores
it and \code{encode} writes it as zeros.

\begin{display}
\renewcommand{\arraystretch}{1.25}
\begin{tabular}{@{}l@{\hspace{6pt}}*{16}{@{}>{\centering\arraybackslash}p{\bitwidth}@{}}%
  >{\hspace{10pt}}l@{}}
 & \bitnumber{15} & \bitnumber{14} & \bitnumber{13} & \bitnumber{12} & \bitnumber{11} &
   \bitnumber{10} & \bitnumber{9} & \bitnumber{8} & \bitnumber{7} & \bitnumber{6} &
   \bitnumber{5} & \bitnumber{4} & \bitnumber{3} & \bitnumber{2} & \bitnumber{1} &
   \bitnumber{0} & \\
\cline{2-17}
\label{fmt:F0}F0 & \fld{5}{opcode} & \fldlast{11}{\code{00000000000}} &
  \small NOP, HLT, RET \\ \cline{2-17}
\label{fmt:F1}F1 & \fld{5}{opcode} & \fld{3}{\code{rd}} & \fld{3}{\code{rs1}} &
  \fld{3}{\code{rs2}} & \fldlast{2}{\code{00}} &
  \small ADD, SUB, MUL, AND, OR, XOR, SLT \\ \cline{2-17}
\label{fmt:F2}F2 & \fld{5}{opcode} & \fld{3}{\code{rd}} & \fld{3}{\code{rs}} &
  \fldlast{5}{\code{00000}} & \small MOV, NOT \\ \cline{2-17}
\label{fmt:F3}F3 & \fld{5}{opcode} & \fld{3}{\code{rd}} &
  \fldlast{8}{\code{imm} or \code{shamt}} & \small LIL, LIH, ADDI, SHL, SHR \\ \cline{2-17}
\label{fmt:F4}F4 & \fld{5}{opcode} & \fld{3}{\code{rd}, \code{rs}} & \fld{3}{\code{base}} &
  \fldlast{5}{\code{off}} & \small LDW, STW, LDB, STB \\ \cline{2-17}
\label{fmt:F5}F5 & \fld{5}{opcode} & \fld{3}{\code{rs}} & \fldlast{8}{\code{off}} &
  \small BZ, BN \\ \cline{2-17}
\label{fmt:F6}F6 & \fld{5}{opcode} & \fldlast{11}{\code{off}} & \small JMP, CALL \\ \cline{2-17}
\label{fmt:F7}F7 & \fld{5}{opcode} & \fld{3}{\code{rs}, \code{rd}} &
  \fldlast{8}{\code{00000000}} & \small PUSH, POP \\ \cline{2-17}
\end{tabular}
\end{display}

\subsection{Opcodes}\label{sec:opcodes}

Opcodes 0 to 26 are assigned; 27 to 31 are not. The table summarises the \code{decode} clauses
of \autoref{sec:isa}; where it and a clause differ, the clause is right.

\begin{display}
\small
\begin{tabular}{@{}rcc@{\hspace{1.2em}}l@{\hspace{1.2em}}l@{\hspace{1.2em}}cr@{}}
\toprule
Opcode & Hex & Bits & Mnemonic & Syntax & Format & Page \\
\midrule
\forlistloop{\oprow}{\oplist}
27--31 & \code{1B}--\code{1F} & \code{11011}--\code{11111} & \multicolumn{4}{l}{unassigned
  (\code{decode} returns \code{None}, \autoref{sec:codec})} \\
\bottomrule
\end{tabular}
\end{display}

For example, the word \code{0x4C4C} is \code{01001} \code{100} \code{010} \code{011} \code{00}:
opcode 9, format F1, with \code{rd} = 4, \code{rs1} = 2 and \code{rs2} = 3, which is
\code{ADD R4, R2, R3}.

\subsection{Operands}\label{sec:operands}

The syntax lines of \autoref{sec:isa} write \code{Rd}, \code{Rs}, \code{Rs1}, \code{Rs2} and
\code{Rb} for the register operands \code{rd}, \code{rs}, \code{rs1}, \code{rs2} and
\code{base} of the model, and \code{imm}, \code{shamt} and \code{off} for the immediates.

\begin{display}
\small
\begin{tabular}{@{}>{\raggedright\arraybackslash}p{0.27\linewidth}%
  >{\centering\arraybackslash}p{2.2em}%
  p{\dimexpr0.73\linewidth-2.2em-4\tabcolsep\relax}@{}}
\toprule
Field & Bits & Meaning \\
\midrule
\code{rd}, \code{rs}, \code{rs1}, \code{rs2}, \code{base} & 3 & A register number, 0 to 7.
  \code{rd} is the destination; \code{rs}, \code{rs1} and \code{rs2} are sources; \code{base}
  holds a memory address. \\
\code{imm} & 8 & An immediate byte: unsigned for LIL and LIH, two's complement
  ($-128$ to $127$) for ADDI. \\
\code{shamt} & 8 & A shift count; only bits 4 to 0 are used. \\
\code{off} (LDW, STW, LDB, STB) & 5 & A signed offset in bytes, $-16$ to $15$. \\
\code{off} (BZ, BN) & 8 & A signed offset in instructions, $-128$ to $127$. \\
\code{off} (JMP, CALL) & 11 & A signed offset in instructions, $-1024$ to $1023$. \\
\bottomrule
\end{tabular}
\end{display}

\subsection{Encoding, decoding and execution}\label{sec:codec}

The type \code{instruction} is a union with one constructor per mnemonic; the constructors are
declared by the \code{union clause} lines of \code{tara.sail}, with the operand types of the
fields in the \code{decode} clauses of \autoref{sec:isa}. Three scattered functions are defined
on it, one clause per instruction:
\begin{itemize}
\item \sym{tarazencode}{encode} turns an instruction into its 16-bit word: the opcode, then the
  fields, with the padding written as zeros;
\item \sym{tarazdecode}{decode} turns a word into an instruction, or \code{None}. Each clause
  matches the opcode in bits 15 to 11 and ignores the padding, so every word with an assigned
  opcode decodes and several words give the same instruction;
\item \sym{tarazexecute}{execute} carries out an instruction. It reads and writes the state of
  \autoref{sec:state}; it sees PC as the address of the instruction itself and proposes the
  successor in \sym{tararegisterznextPC}{nextPC}, which the drivers commit
  (\autoref{sec:retire}). An instruction reads its sources before it writes its result;
  \mnem{PUSH} and \mnem{POP} are the only ones with two steps, and they perform them in the order
  of their clauses.
\end{itemize}

\keep{\taratypeinstruction\taravalencode\taravaldecode\taravalexecute}

The opcodes 27 to 31 are unassigned; the last clause of \code{decode} gives them no
instruction.

\keep{\tarafclNonedecode}

% =============================================================================================
\section{The instruction set}\label{sec:isa}
% =============================================================================================

Each entry gives the syntax, the opcode and the format of an instruction, a short description
and the clauses of \code{encode}, \code{decode} and \code{execute} that define it. An
instruction retires when its effect has been applied and PC has moved on (\autoref{sec:retire}).

% ---------------------------------------------------------------------------------------------
\group{System and data movement}{sec:isa-system}
% ---------------------------------------------------------------------------------------------

\mnem{NOP} and \mnem{HLT} take no operands. \mnem{MOV} copies a register; \mnem{LIL} and
\mnem{LIH} build a 16-bit constant one byte at a time.

\begin{instruction}{NOP}{Does nothing. The instruction retires and PC moves to the next word.}
\tarafclNOPencode
\tarafclNOPdecode
\tarafclNOPexecute
\end{instruction}

\begin{instruction}{HLT}{Stops the CPU. HLT retires like any instruction, so PC advances by 2,
  and then sets the halt latch \sym{tararegisterzHALTED}{HALTED}. From then on
  \sym{tarazstep}{step} and \sym{tarazrunzyinstruction}{run\_instruction} return
  \sym{tarazStopped}{Stopped} until \sym{tarazreset}{reset} clears the latch
  (\autoref{sec:halting}).}
\tarafclHLTencode
\tarafclHLTdecode
\tarafclHLTexecute
\end{instruction}

\begin{instruction}{MOV}{Copies \code{Rs} to \code{Rd}.}
\tarafclMOVencode
\tarafclMOVdecode
\tarafclMOVexecute
\end{instruction}

\begin{instruction}{LIL}{Loads the immediate into the low byte of \code{Rd} and clears the upper
  byte.}
\tarafclLILencode
\tarafclLILdecode
\tarafclLILexecute
\end{instruction}

\begin{instruction}{LIH}{Loads the immediate into the upper byte of \code{Rd} and leaves the
  low byte unchanged. LIL followed by LIH builds any 16-bit constant.}
\tarafclLIHencode
\tarafclLIHdecode
\tarafclLIHexecute
\end{instruction}

% ---------------------------------------------------------------------------------------------
\group{Memory}{sec:isa-memory}
% ---------------------------------------------------------------------------------------------

A memory instruction names a data register, a base register and a signed 5-bit offset in bytes.
The effective address is the sum, computed modulo $2^{16}$; the rules of \autoref{sec:memory}
then apply to it: only its low 11 bits count, and a word access ignores bit 0. Words are
big-endian, and a read of \code{0x5FF} returns the input port (\autoref{sec:input}).

\begin{instruction}{LDW}{Loads the word at the effective address \code{Rb + off} into
  \code{Rd}.}
\tarafclLDWencode
\tarafclLDWdecode
\tarafclLDWexecute
\end{instruction}

\begin{instruction}{STW}{Stores \code{Rs} as the word at the effective address
  \code{Rb + off}, high byte first.}
\tarafclSTWencode
\tarafclSTWdecode
\tarafclSTWexecute
\end{instruction}

\begin{instruction}{LDB}{Loads the byte at the effective address \code{Rb + off} into the low
  byte of \code{Rd} and clears the upper byte.}
\tarafclLDBencode
\tarafclLDBdecode
\tarafclLDBexecute
\end{instruction}

\begin{instruction}{STB}{Stores the low byte of \code{Rs} at the effective address
  \code{Rb + off}.}
\tarafclSTBencode
\tarafclSTBdecode
\tarafclSTBexecute
\end{instruction}

% ---------------------------------------------------------------------------------------------
\group{Arithmetic}{sec:isa-arithmetic}
% ---------------------------------------------------------------------------------------------

Arithmetic is modulo $2^{16}$. There are no carry or overflow flags.

\begin{instruction}{ADD}{Adds \code{Rs1} and \code{Rs2}, modulo $2^{16}$.}
\tarafclADDencode
\tarafclADDdecode
\tarafclADDexecute
\end{instruction}

\begin{instruction}{SUB}{Subtracts \code{Rs2} from \code{Rs1}, modulo $2^{16}$.}
\tarafclSUBencode
\tarafclSUBdecode
\tarafclSUBexecute
\end{instruction}

\begin{instruction}{ADDI}{Adds the immediate, sign-extended from 8 bits ($-128$ to $127$), to
  \code{Rd}, modulo $2^{16}$.}
\tarafclADDIencode
\tarafclADDIdecode
\tarafclADDIexecute
\end{instruction}

\begin{instruction}{MUL}{Stores the low 16 bits of the product of \code{Rs1} and \code{Rs2}.
  The low 16 bits are the same whether the operands are read as signed or as unsigned.}
\tarafclMULencode
\tarafclMULdecode
\tarafclMULexecute
\end{instruction}

% ---------------------------------------------------------------------------------------------
\group{Logic}{sec:isa-logic}
% ---------------------------------------------------------------------------------------------

The logic instructions work on all 16 bits at once.

\begin{instruction}{AND}{Bitwise and of \code{Rs1} and \code{Rs2}.}
\tarafclANDencode
\tarafclANDdecode
\tarafclANDexecute
\end{instruction}

\begin{instruction}{OR}{Bitwise or of \code{Rs1} and \code{Rs2}.}
\tarafclORencode
\tarafclORdecode
\tarafclORexecute
\end{instruction}

\begin{instruction}{XOR}{Bitwise exclusive or of \code{Rs1} and \code{Rs2}.}
\tarafclXORencode
\tarafclXORdecode
\tarafclXORexecute
\end{instruction}

\begin{instruction}{NOT}{Bitwise complement of \code{Rs}.}
\tarafclNOTencode
\tarafclNOTdecode
\tarafclNOTexecute
\end{instruction}

% ---------------------------------------------------------------------------------------------
\group{Shifts}{sec:isa-shifts}
% ---------------------------------------------------------------------------------------------

Both shifts are logical and take their count from the low five bits of the immediate
(``\code{shamt = imm8[4:0]}''); bits 7 to 5 of the field are ignored. The RTL does not say
what a count of 16 or more does; the model shifts every bit out, giving zero
(\autoref{app:assumptions}).

\begin{instruction}{SHL}{Shifts \code{Rd} left by \code{shamt[4:0]}; zeros enter at the
  right.}
\tarafclSHLencode
\tarafclSHLdecode
\tarafclSHLexecute
\end{instruction}

\begin{instruction}{SHR}{Shifts \code{Rd} right by \code{shamt[4:0]}; zeros enter at the
  left.}
\tarafclSHRencode
\tarafclSHRdecode
\tarafclSHRexecute
\end{instruction}

% ---------------------------------------------------------------------------------------------
\group{Comparison}{sec:isa-comparison}
% ---------------------------------------------------------------------------------------------

\begin{instruction}{SLT}{Sets \code{Rd} to 1 if \code{Rs1} is less than \code{Rs2}, both read
  as signed 16-bit numbers, and to 0 otherwise.}
\tarafclSLTencode
\tarafclSLTdecode
\tarafclSLTexecute
\end{instruction}

% ---------------------------------------------------------------------------------------------
\group{Control flow}{sec:isa-control}
% ---------------------------------------------------------------------------------------------

A relative target is $\mathrm{PC} + 2 + 2\cdot\mathrm{off}$: the offset counts instructions
and is relative to the instruction that follows. PC keeps only 11 bits, so targets wrap around
the 2~KB memory (\autoref{sec:retire}). BZ and BN test their register itself, never the status
flags.

\begin{instruction}{BZ}{Branches if \code{Rs} is zero. Otherwise execution continues with the
  next instruction.}
\tarafclBZencode
\tarafclBZdecode
\tarafclBZexecute
\end{instruction}

\begin{instruction}{BN}{Branches if \code{Rs} is negative, that is, if its bit 15 is set.
  Otherwise execution continues with the next instruction.}
\tarafclBNencode
\tarafclBNdecode
\tarafclBNexecute
\end{instruction}

\begin{instruction}{JMP}{Jumps. The 11-bit offset ($-1024$ to $1023$ instructions) spans the
  whole memory.}
\tarafclJMPencode
\tarafclJMPoffdecode
\tarafclJMPexecute
\end{instruction}

\begin{instruction}{CALL}{Calls a subroutine: writes the address of the next instruction,
  masked to 11 bits, to the link register R6 and jumps like JMP. CALL touches neither R7 nor
  memory.}
\tarafclCALLencode
\tarafclCALLoffdecode
\tarafclCALLexecute
\end{instruction}

\begin{instruction}{RET}{Returns: PC becomes the value of the link register R6, of which PC
  keeps the low 11 bits.}
\tarafclRETencode
\tarafclRETdecode
\tarafclRETexecute
\end{instruction}

% ---------------------------------------------------------------------------------------------
\group{Stack}{sec:isa-stack}
% ---------------------------------------------------------------------------------------------

The stack is the memory addressed by R7. It grows downwards, one word at a time; there is no
overflow check, and a program sets R7 before the first PUSH, because reset leaves registers
alone. The two instructions follow the order of the ISA reference table, which matters when the
register is R7 itself (\autoref{app:assumptions}).

\begin{instruction}{PUSH}{Decrements R7 by 2, then stores \code{Rs} at the address in R7.
  \code{Rs} is read after the decrement, so PUSH R7 stores the decremented stack pointer.}
\tarafclPUSHencode
\tarafclPUSHrsdecode
\tarafclPUSHexecute
\end{instruction}

\begin{instruction}{POP}{Loads the word at the address in R7 into \code{Rd}, then increments R7
  by 2. The increment comes last, so POP R7 leaves the loaded value plus 2.}
\tarafclPOPencode
\tarafclPOPrddecode
\tarafclPOPexecute
\end{instruction}

% =============================================================================================
\section{Execution}\label{sec:execution}
% =============================================================================================

The drivers of \code{step.sail} run one instruction per call. \sym{tarazstep}{step} and
\sym{tarazrunzyinstruction}{run\_instruction} take the five input lines as their first argument
and hold them for the duration of the call.

\subsection{Results}\label{sec:results}

A driver reports what happened with a \sym{taratypezretirement}{retirement}.

\keep{\taratyperetirement}

\begin{itemize}
\item \phantomsection\label{tarazRetired}\code{Retired}: an instruction ran to completion.
  HLT retires like any other instruction; it also sets the halt latch.
\item \phantomsection\label{tarazStopped}\code{Stopped}: the CPU is halted and nothing
  happened: no fetch, no input latched, no state changed.
\item \phantomsection\label{tarazIllegal}\code{Illegal}: the fetched word has an unassigned
  opcode (27 to 31); the payload is that word. The fetch advanced PC past the word; registers,
  memory and the halt latch are unchanged. The model does not halt the CPU: what the hardware
  does next is not published, so the host decides.
\end{itemize}

\subsection{Retiring an instruction}\label{sec:retire}

\sym{tarazretire}{retire} runs a decoded instruction. The default successor is PC + 2;
\sym{tarazexecute}{execute} may replace it in \sym{tararegisterznextPC}{nextPC}, and the result,
masked to 11 bits, becomes PC. An instruction therefore sees PC as its own address, while the RTL
has already incremented PC by the time the instruction executes; the PC-relative targets of
\autoref{sec:isa} are written accordingly.

\keep{\taravalretire\tarafnretire}

\subsection{Fetch, decode, execute}\label{sec:step}

\sym{tarazstep}{step} is the model of the processor: it runs the next instruction from memory.
\begin{enumerate}
\item If \sym{tararegisterzHALTED}{HALTED} is set it returns \sym{tarazStopped}{Stopped}.
\item It latches the input lines in \sym{tararegisterzKEYS}{KEYS}.
\item It fetches the word at PC with \sym{tarazreadzyword}{read\_word}. The fetch is an ordinary
  word read:
  it ignores bit 0 of PC, and it sees the input port if it covers \code{0x5FF}. It also sees
  the stores of earlier instructions, which is how self-modifying code works.
\item It decodes the word. For an assigned opcode it retires the instruction.
\item For an unassigned opcode it advances PC to \sym{tarazpczymask}{pc\_mask}\code{(PC + 2)},
  as the fetch steps of the RTL do for every word, and returns \sym{tarazIllegal}{Illegal} with
  the word.
\end{enumerate}

\keep{\taravalstep\tarafnstep}

\subsection{Running a decoded instruction}\label{sec:run}

\sym{tarazrunzyinstruction}{run\_instruction} executes an instruction that is already decoded,
as if it had been fetched from PC. It does no fetch; a halted CPU does nothing, as for
\sym{tarazstep}{step}. It runs an instruction without placing it in memory.

\keep{\taravalrunInstruction\tarafnrunInstruction}

\subsection{Halting and reset}\label{sec:halting}

\mnem{HLT} sets \sym{tararegisterzHALTED}{HALTED}. Once it is set, \sym{tarazstep}{step} and
\sym{tarazrunzyinstruction}{run\_instruction} return \sym{tarazStopped}{Stopped} without
touching the state, until \sym{tarazreset}{reset}, which ``clears halted state, and restarts from
PC 0''; registers and memory are unchanged.

\keep{\taravalreset\tarafnreset}

% =============================================================================================
\section{Assembly syntax}\label{sec:syntax}
% =============================================================================================

\code{syntax.sail} prints an instruction in the syntax of the TARA Studio assembler. Registers
print as \code{R0} to \code{R7}. The immediates of LIL, LIH, SHL and SHR print as unsigned
decimal numbers; the immediate of ADDI and all offsets print signed. A memory operand prints as
\code{off(Rb)}. Branch and jump targets print as raw signed instruction offsets, so
assembling the text gives back the canonical encoding. The Studio assembler also accepts a
label in place of an offset. The file overloads \code{\textasciicircum} to concatenate strings.

\keep{\taraoverloadCzEightoperatorzZerozQzNine\taravalreg\tarafnreg\taravalmem\tarafnmem}

\keep{\taravalopOne\tarafnopOne\taravalopTwo\tarafnopTwo\taravalopThree\tarafnopThree}

% The longest listing may continue on the next page.
\par\nopagebreak\vspace{3pt}
{\setlength{\parskip}{0pt}\taravalassembly\tarafnassembly}

\clearpage
\appendix
% =============================================================================================
\section{Assumptions}\label{app:assumptions}
% =============================================================================================

Where the RTL description is silent the model assumes the following.
\begin{itemize}
\item Shift counts of 16 to 31 (``\code{shamt = imm8[4:0]}'') shift out every bit
  (\autoref{sec:isa-shifts}).
\item PUSH R7 stores the decremented stack pointer and POP R7 leaves the loaded value plus 2.
  This follows the order of the ISA reference table, and TARA Studio agrees.
\item A word read that covers \code{0x5FF}, by LDW or by a fetch at \code{0x5FE}, gets the
  input byte as its low half (\autoref{sec:memory}).
\item A store to \code{0x5FF} writes the memory byte, which reads never return. Memory has a
  single write port, and a byte store is a read-modify-write of a word.
\item After the fetch of an opcode from 27 to 31 the model reports \code{Illegal} and does not
  halt; the control ROM that would decide what happens next is unpublished
  (\autoref{sec:results}).
\item Reset clears the halt latch and sets PC to 0; registers and memory keep their values
  (\autoref{sec:halting}).
\item The input lines hold still for the duration of one instruction
  (\autoref{sec:input}).
\end{itemize}

% =============================================================================================
\section{Differences from TARA Studio}\label{app:studio}
% =============================================================================================

Where the RTL description and the TARA Studio simulator (taracpu 1.2.2) disagree, the model
follows the RTL description.

\begin{display}
\small
\begin{tabular}{@{}>{\raggedright\arraybackslash}p{0.28\linewidth}%
  >{\raggedright\arraybackslash}p{0.34\linewidth}%
  >{\raggedright\arraybackslash}p{\dimexpr0.38\linewidth-4\tabcolsep\relax}@{}}
\toprule
 & RTL description (the model) & TARA Studio \\
\midrule
Word access or fetch at an odd address & bit 0 is ignored: ``word index =
  addr \guillemotright{} 1''
  & bytes \code{addr} and \code{addr+1} \\
Word access or fetch at \code{0x7FF} & the word at \code{0x7FE} & error; the CPU halts \\
CALL at \code{0x7FE} & links \code{0x000}: PC is ``masked to 0x7FF'' & links \code{0x800} \\
Reading \code{0x5FF} & the live input lines & memory, rewritten by the UI on key events \\
Opcodes 27 to 31 & fetched like any word, so PC + 2 & error; the CPU halts, PC unchanged \\
Q key & requests a reset; QUIT is the switch sw[4] & sets QUIT (bit 4 of \code{0x5FF}) \\
\bottomrule
\end{tabular}
\end{display}

% =============================================================================================
\section{Example: Fibonacci}\label{app:example}
% =============================================================================================

The Fibonacci program of the hardware manual shows encodings and relative targets at work. R2
and R3 hold two consecutive Fibonacci numbers and R5 counts the ten iterations. The assembly
column is what \sym{tarazassembly}{assembly} prints.

\begin{display}
\small
\begin{tabular}{@{}lll>{\raggedright\arraybackslash}p{0.4\linewidth}@{}}
\toprule
Address & Word & Assembly & Target \\
\midrule
\code{0x000} & \code{0x1A00} & \code{LIL R2, 0} & \\
\code{0x002} & \code{0x1B01} & \code{LIL R3, 1} & \\
\code{0x004} & \code{0x1D0A} & \code{LIL R5, 10} & \\
\code{0x006} & \code{0xA505} & \code{BZ R5, 5} &
  if R5 is 0: $\mathtt{0x006} + 2 + 2\cdot5 = \mathtt{0x012}$ \\
\code{0x008} & \code{0x4C4C} & \code{ADD R4, R2, R3} & \\
\code{0x00A} & \code{0x1260} & \code{MOV R2, R3} & \\
\code{0x00C} & \code{0x1380} & \code{MOV R3, R4} & \\
\code{0x00E} & \code{0x5DFF} & \code{ADDI R5, -1} & \\
\code{0x010} & \code{0xB7FA} & \code{JMP -6} &
  $\mathtt{0x010} + 2 - 2\cdot6 = \mathtt{0x006}$ \\
\code{0x012} & \code{0x1640} & \code{MOV R6, R2} & \\
\code{0x014} & \code{0x0800} & \code{HLT} & \\
\bottomrule
\end{tabular}
\end{display}

Started at PC 0, the program retires 66 instructions and halts with R2 = \code{0x0037} (55, the
tenth Fibonacci number), R3 = R4 = \code{0x0059}, R5 = 0 and R6 = \code{0x0037}. PC is
\code{0x0016}: HLT retires, so PC has advanced past it.

\end{document}
